% F7: Literature review. 1.5 to 2 pages.

%   - Organized by theme or approach, not paper by paper.
%   - IEEE citations throughout.
%   - Ends in an explicit gap tied to Heilmeier question 2.
% Before writing: the synthesis matrix (6 to 8 papers, every cell checked and
% marked verified), one map, and the argument outline go in
% appendices/my_material.tex. The map goes in figures/.
% If AI drafts this section from your outline and matrix: save its draft
% unchanged as sections/literature_review_ai_v0.tex, then edit this file, and
% add \trackededits{literature_review_ai_v0} to appendices/tracked_edits.tex.
% Winter (W6): this review is condensed into your paper's related-work section.

\section{Literature Review}

Integrating a water-based resistojet onto a resource-constrained CubeSat requires coordinating electrical power allocation, attitude slewing, and state estimation across changing orbital conditions. Unlike larger spacecraft whose buses power thrusters directly, a CubeSat using a phase-change material (PCM) thermal battery must accumulate thermal energy gradually before executing an orbital burn. The goal is to minimize the time required to reach a burn while maintaining the pointing requirements associated with firing. In addition, some form of backup tumbling estimation in the post-burn period is useful to determine whether another firing would be dangerous. Reviewing existing literature across power scheduling, attitude optimization, and attitude estimation reveals both foundational methods and critical gaps for low-power propulsive CubeSats.

This mission is specifically well-suited for power scheduling because its goal is to maximize available power over various load conditions up to the propulsion burn. The propulsion system must be fed all available power to phase-change the material in as little time as possible. From experiments, reaching operating temperature takes at least two hours at the maximum heater power of 40 W, on the order of a single orbit in LEO. Given that CubeSat power budgets are well under the maximum available heater power \cite{clarkHugePowerDemandItsyBitsy2010}, keeping the satellite generally operational over multiple orbits while charging is critical.

Several works focus on optimizing satellite loads over long horizons, attempting to maximize available power, battery life, or stored energy. Essentially all formulations of power delivery systems adapt which loads receive power through load scheduling. Communication-focused CubeSats may benefit from fixed-horizon scheduling when the load and transmission conditions are well understood. In the case of communications, favorable channel conditions may make it worthwhile to drive transmit power up during certain periods in order to minimize the bit error rate (BER) \cite{husseinOptimizedLoadSchedulingAlgorithm2023}. Contrast this with load conditions that are unknown in the future. A self-adaptive model-predictive control (MPC) algorithm can instead manage multiple loads, taking into account future power generation and loading conditions in different parts of the orbit \cite{akarneSelfAdaptiveEnergyAwarePower2026}. Part of the unique value of the MPC method is that it adapts to varying CubeSat mission goals. For example, MPC parameters can be tuned over an orbit to optimize for energy storage or battery state of health (SOH) \cite{akarneSelfAdaptiveEnergyAwarePower2026}. Part of the challenge of load scheduling is the computation required to estimate the best load schedule beforehand. Other work focuses on computational speed across coupled subsystems, implementing a fully convex formulation of a similar problem \cite{doff-sottaSpacecraftEnergyManagement2022}. The drawback to this formulation is its limited ability to adapt to changing conditions.

In the case of this mission, much of this work falls short of the needs of an average CubeSat and the goals of the propulsion mission. Work that optimizes for a specific payload tends to lump all non-payload power into a ``fixed-load'' category \cite{husseinOptimizedLoadSchedulingAlgorithm2023}. Any method that does not predict future loading may suffer when increased demand forces increased draw. Some power management schemes do not even use feedback in the control loop, in order to minimize the complexity of the power management method. Predictions of power generation and load availability are instead made from an estimated orbit, assuming no errors in the model \cite{doff-sottaSpacecraftEnergyManagement2022}. Maintaining a steady supply of energy even through eclipse conditions is important here, as having to shut off the heating element during every eclipse would increase charge times. Methods which iteratively predict power usage and generation over time are much more successful at maintaining a load over a long period of time \cite{akarneSelfAdaptiveEnergyAwarePower2026}.

Turning to the thruster itself, any burn will be aligned with a single spacecraft axis, so communication, solar charging, and firing are potentially at odds. The goal for this mission is to maximize power generation during multiple orbits while meeting the pointing goals required by a single-axis thruster.

Optimal slewing has been demonstrated for both optimal-time and optimal-power formulations. For a single-thruster spacecraft maneuvering across non-concentric orbits, some work focuses on solving for burn locations and $\Delta v$ requirements, which then drive pointing decisions. One method utilizes a heuristic optimization (particle swarm optimization) that reduces the optimal-time pointing formulation to a comparatively simple problem \cite{zhangFinitetimeAttitudeSetpoint2020}. Maximizing power during slewing is significantly more complex, as a solver must take into account pointing goals that are at odds with each other while minimizing control action. The result is a nonlinear optimization formulation that generates $\approx 40\%$ more energy during the maneuver than a PD controller \cite{kristiansenEnergyOptimalAttitude2024}\cite{kristiansenEnergyOptimalAttitude2021}.

The last key to ensuring propulsive burn reliability is post-burn tumbling estimation. Gyroscopes and three-axis magnetometers (TAM) are nearly universal on CubeSats, and an unscented Kalman filter (UKF) or particle filter (PF) can use their measurements for attitude state estimation over an orbit. For large initial errors, it has been shown that the PF converges faster than a UKF on the same measurements, at the cost of computational complexity \cite{xue-yuanjiangSpacecraftAttitudeEstimation2005}. This motivates a sequential PF-to-UKF algorithm to estimate tumbling after firing while reducing computation time.

The limit of current practice (Heilmeier Q2) is that a CubeSat resistojet draws more instantaneous power than a 3U bus generates. A PCM thermal battery removes that constraint by decoupling burn power from charge power (Heilmeier Q3), but it replaces an instantaneous power problem with a scheduling problem no reviewed work addresses. Multi-orbit, high-power charge events will require some scheme to handle the strain placed on the system and to leave residual power for the recovery. Scheduling work \cite{akarneSelfAdaptiveEnergyAwarePower2026}\cite{doff-sottaSpacecraftEnergyManagement2022}\cite{husseinOptimizedLoadSchedulingAlgorithm2023} assumes fixed orbital geometry or does not carry the high-load condition of this system. Pointing work \cite{zhangFinitetimeAttitudeSetpoint2020}\cite{kristiansenEnergyOptimalAttitude2024}\cite{kristiansenEnergyOptimalAttitude2021} manages optimal-time transfers or optimal-power scheduling, but carries no thrust objective and is not real-time feasible. Attitude estimation work \cite{xue-yuanjiangSpacecraftAttitudeEstimation2005} proves out PF and UKF methods, but does not compare them against absolute error. Combining these methods and validating their utility to meet the goals of the specific application is the novelty of this work.
